% Ерөнхий дүгнэлт

\phantomsection
\addchaptertocentry{Ерөнхий дүгнэлт}
\label{Summary} % Энэ бүлэг рүү ишлэл хийх бол
%-------------------------------------------------------------------------------
%	Summary
%-------------------------------------------------------------------------------
\section*{Ерөнхий дүгнэлт}

\paragraph{}
Энэхүү бакалаврын төгсөлтийн ажлын хүрээнд ажил горилогчийн ур чадварыг хиймэл оюунаар үнэлэх цогц веб систем боловсруулж, амжилттай хэрэгжүүлсэн. Ажлын зорилгод тавигдсан шаардлагуудыг дараах байдлаар биелүүлсэн.

\paragraph{}
\textbf{Онол, арга зүйн судалгааны хүрээнд} ажилд авах уламжлалт болон AI-д суурилсан аргуудыг харьцуулан шинжилж, CV болон ажлын байрны нийцлийг тодорхойлоход ашиглагддаг машин сургалт (SVM, KNN), байгалийн хэл боловсруулалт (NER, Sentence-BERT), Siamese network болон томоохон хэлний загвар (LLM) зэрэг алгоритмуудыг судалсан. Монголын зах зээлд үйл ажиллагаа явуулдаг zangia.mn болон lambda.global системүүдтэй харьцуулалт хийсний үндсэн дээр тайлбарлах боломжтой (explainable) үнэлгээ гаргах давуу талтай LLM-д суурилсан аргыг сонгосон. Мөн системийн хөгжүүлэлтэд ашиглах Go, Nuxt 4, PostgreSQL, Anthropic Claude API зэрэг технологиудыг тодорхойлсон.

\paragraph{}
\textbf{Системийн шаардлага болон шинжилгээний хүрээнд} ажил горилогч, ажил олгогч, зочин гэсэн гурван тоглогчтой 12 юзкейс бүхий юзкейс диаграммыг байгуулж, функциональ болон функциональ бус шаардлагуудыг тодорхойлсон. Шинжилгээний шатны класс, дараалал болон үйл ажиллагааны диаграммуудыг гаргаж, системийн зохион байгуулалт болон процессын дарааллыг тодорхой болгосон.

\paragraph{}
\textbf{Системийн зохиомж болон хөгжүүлэлтийн хүрээнд} модульчилсан монолит архитектурыг сонгон, back-end-ийг \texttt{aiman}, \texttt{jobman}, \texttt{appman}, \texttt{taskman}, \texttt{cvman}, \texttt{userman} зэрэг домэйнчилсан модулиудад хувааж Go хэл дээр хөгжүүлсэн. Claude Sonnet загварт тулгуурласан CV үнэлгээний \texttt{AssessCV}, даалгавар автоматаар үүсгэх \texttt{GenerateTask} болон CV мэдээлэл задлах \texttt{ParseCV} функцуудыг хэрэгжүүлсэн. Ярилцлагын цаг товлоход Google Calendar API-тай интеграц хийж, AI үнэлгээний үр дүнг WebSocket дамжуулан бодит цагт хэрэглэгчид хүргэдэг болгосон. Хэрэглэгчийн интерфэйсийг Nuxt 4, Tailwind CSS болон shadcn/vue ашиглан хөгжүүлсэн.

\paragraph{}
\textbf{Төслийн менежментийн хувьд} хөгжүүлэлтийн явцад Git хувилбарын удирдлагын систем болон GitHub платформыг ашиглан кодын өөрчлөлтийг хянаж, хамтран ажиллах боломжийг бүрдүүлсэн. Кодын санг \texttt{github.com} дээр нийтийн бус \textit{(private)} хэлбэрт байршуулсан. Систем нь одоогоор локал орчинд ажиллаж байгаа бөгөөд \texttt{make web} болон \texttt{make ui} командуудаар хурдан ажиллуулах боломжтой.

\paragraph{}
\textbf{Цаашдын хөгжүүлэлтийн талаар} системийг Docker контейнержуулалт ашиглан байгууллагын серверт байршуулах, Nginx reverse proxy-оор front-end болон back-end серверийг нэгтгэх боломжтой. Мөн дараах чиглэлүүдэд системийг цаашид хөгжүүлэх боломжтой:

\begin{itemize}
    \item \textbf{CV файлаас профайл автоматаар бөглөх:} Одоогийн хэрэгжилтэд \texttt{ParseCV} функц нь CV-ийн мэдээллийг задлах боломжтой боловч цаашид ажил горилогч CV файл оруулах бүрт системд бүртгэлтэй профайл нь автоматаар шинэчлэгдэх болгон, давхардсан мэдээлэл оруулах хэрэгцээг арилгах;
    \item \textbf{Push мэдэгдлийн систем:} Ажил горилогч болон ажил олгогчид шинэ ажлын зар, анкетын статус өөрчлөлт, даалгаврын хугацааны сануулга зэргийг имейл болон browser push notification-оор хүргэх;
    \item \textbf{Тохирох ажлын байр (Recommendation):} Ажил горилогчийн CV профайл дахь ур чадвар, туршлага, боловсрол болон өмнөх үнэлгээний үр дүнд тулгуурлан тухайн хэрэглэгчид хамгийн тохирох ажлын байруудыг автоматаар санал болгох;
    \item \textbf{Чатботын хэсэгт RAG нэмэх:} Одоогийн чатботын хэрэгжилтэд ажлын байрны мэдээллийг prompt-д шууд дамжуулдаг бол цаашид Retrieval-Augmented Generation (RAG) аргыг ашиглан ажлын зарын мэдээллийн сан, байгууллагын дотоод баримт бичиг, нийтлэгдсэн мэдээлэл зэргийг векторжуулан хадгалж, хэрэглэгчийн асуултад хамаарах хэсгийг динамикаар авч prompt-д нэмэх замаар чатботын хариултын нарийвчлал болон мэдлэгийн хүрээг ихэсгэх.
\end{itemize}

\paragraph{}
Нийтдээ энэхүү ажлаар ажил горилогчийн CV-г автоматаар шүүх, тайлбарлах боломжтой AI үнэлгээ хийх, ур чадварын практик даалгавар илгээх болон ярилцлага товлох зэрэг сонгон шалгаруулалтын бүхий л үе шатыг нэг платформ дотор хөгжүүлснээр уламжлалт ажилд авах процессын цаг хугацааны болон субъектив алдааны асуудлыг шийдвэрлэх боломжийг бүрдүүлж, тавьсан зорилгод хүрсэн.